\section{Nonlinear sources and the limits of linear trapping}
\label{nl:section}

A localized eigenmode of the linearized equation need not remain
radiationless at finite amplitude. We distinguish an exact source
calculation, the kinematic opening of radiation channels, and a nonzero
outgoing response. Existing nonlinear evidence for radial location C0
is described below; it is not transferred to the additional modes T1
and T2 without evaluating their own sources.

\subsection{Quadratic expansion}
\label{nl:expansion}
We use the complex nonlinear Klein--Gordon equation
\begin{equation}
 \phi_{tt}-\Delta\phi+U'(|\phi|^2)\phi=0,
 \qquad U(S)=S-S^2+\tfrac12 S^3,
 \qquad \phi=e^{i\omega t}(f+z),
 \label{nl:model}
\end{equation}
where the real background satisfies
$-\Delta f-\omega^2f+U'(f^2)f=0$. Setting $s=f^2$, the linear operator is
\begin{equation}
 \mathcal Lz=z_{tt}+2i\omega z_t-\Delta z
 +[U'(s)+sU''(s)-\omega^2]z+sU''(s)\overline z.
 \label{nl:linear}
\end{equation}
The quadratic term in the equation's left-hand side is exactly
\begin{align}
 G_2[z]&=fU''(s)(z^2+2|z|^2)
       +\tfrac12 f^3U'''(s)(z+\overline z)^2 \nonumber\\
       &=\mathcal A z^2+2\mathcal A|z|^2+\mathcal C\overline z^{\,2},
 &\mathcal A&=f\left(-2+\tfrac92f^2\right),
 \qquad \mathcal C=\tfrac32f^3.
 \label{nl:quadratic}
\end{align}
Consequently $z=\epsilon z_1+\epsilon^2z_2+\cdots$ gives
$\mathcal Lz_1=0$ and $\mathcal Lz_2=-G_2[z_1]$. In particular, the
forcing has the opposite sign to the Taylor term in
Eq.~\eqref{nl:quadratic}.

For a real spherical harmonic $Y$ and real radial functions $a,b$, write
\begin{equation}
 z_1=Y\bigl(ae^{i\rho t}+be^{-i\rho t}\bigr).
\end{equation}
Then
\begin{align}
 G_2[z_1]&=Y^2\bigl(Q_+e^{2i\rho t}+Q_0+Q_-e^{-2i\rho t}\bigr),\\
 Q_+&=\mathcal A(a^2+2ab)+\mathcal Cb^2,\qquad
 Q_-=\mathcal A(b^2+2ab)+\mathcal Ca^2,\\
 Q_0&=2\mathcal A(a^2+b^2)+2(\mathcal A+\mathcal C)ab.
 \label{nl:coefficients}
\end{align}
These expressions introduce rotating-frame frequencies $0,\pm2\rho$,
or laboratory frequencies $\omega,\omega\pm2\rho$. The zero-frequency
term is static only in the rotating frame.

\subsection{Angular selection and forced channel problems}
\label{nl:selection}
For T1, the radial $l=0$ mode, $Y=Y_{00}$ and
$Y_{00}^2=Y_{00}/\sqrt{4\pi}$, so only $L=0$ occurs at this order.
For T2, choose a real normalized $l=1$ harmonic. After a spatial rotation
it equals $Y_{10}$, and
\begin{equation}
 Y_{10}^2=\frac{Y_{00}}{\sqrt{4\pi}}
          +\frac{Y_{20}}{\sqrt{5\pi}}.
 \label{nl:angular}
\end{equation}
Thus each of the three temporal source coefficients can drive $L=0$
and $L=2$; the antisymmetric $L=1$ product is absent. A different real
orientation rotates the quadrupole within the $L=2$ multiplet. This
standing real mode must not be confused with a rotating complex
$Y_{11}/\overline{Y}_{11}$ block: the latter has
$Y_{11}^2=\sqrt{3/(10\pi)}Y_{22}$ for its corresponding oscillatory product.

Let $c_L$ denote the coefficient in the relevant angular product and
$D_L=-\partial_r^2-2r^{-1}\partial_r+L(L+1)r^{-2}$. For the oscillatory
response $h_{+,L}e^{2i\rho t}+h_{-,L}e^{-2i\rho t}$ one must solve
the coupled system
\begin{equation}
 \begin{pmatrix}
 D_L+U'+sU''-(\omega+2\rho)^2&sU''\\
 sU''&D_L+U'+sU''-(\omega-2\rho)^2
 \end{pmatrix}
 \begin{pmatrix}h_{+,L}\\\overline{h}_{-,L}\end{pmatrix}
 =-c_L\begin{pmatrix}Q_+\\Q_-\end{pmatrix}.
 \label{nl:forced}
\end{equation}
Here $U',U''$ are evaluated at $s$; the conjugated second unknown matters
when imposing outgoing conditions. The real static correction obeys
$[D_L+U'+2sU''-\omega^2]h_{0,L}=-c_LQ_0$.
Gauge and charge constraints must also be fixed in constructing a
nonlinear continuation.

Both T1 and T2 have $|\omega|<1$ and $|\omega\pm2\rho|>1$.
For T2, the latter frequencies are approximately $4.52130$ and $-2.78407$.
The second harmonics are therefore open in the mass-one vacuum, whereas
the rotating-static correction is closed. The centrifugal term does
not change the asymptotic mass threshold. With the convention
$e^{iEt}$, an outgoing wave has radial phase
$e^{-i\operatorname{sgn}(E)\sqrt{E^2-1}\,r}$; a negative $E$ does not
mean negative outward energy transport.

An open channel and a nonzero local source do not establish radiation:
the outgoing source overlap can vanish. A localized continuation without
incoming radiation requires both open amplitudes to vanish in every
allowed $L$ block of Eq.~\eqref{nl:forced}. The T1 response remains open;
an approximate-input T2 calculation is reported below. A nonzero such amplitude would obstruct
a regular localized perturbative continuation with the specified leading
mode, but would not exclude all other nonlinear states. Frequency shifts
of order $\epsilon^2$ do not automatically cancel the second harmonic:
the carrier shift contributes a static background correction at this
order, while a shift of $\rho$ first acts on $\epsilon z_1$ at order
$\epsilon^3$.

\subsection{Complex dipole polarization and a conditional obstruction}
\label{nl:polarization}
The threefold dipole degeneracy allows more than rotations of a real
standing mode. Let $Y_i(\boldsymbol n)=\sqrt{3/(4\pi)}n_i$ be real
orthonormal Cartesian dipole harmonics and write
$A(\boldsymbol n)=\sum_i p_iY_i(\boldsymbol n)$, with
$\boldsymbol p\in\mathbb C^3$. For the same real radial eigenmode as above,
the allowed tangent has the linked sidebands
$z_1=a(r)A e^{i\rho t}+b(r)\overline A e^{-i\rho t}$.
This does not allow independently chosen polarizations of the two
sidebands or a mixture of distinct radial modes.

Set $\nu=\boldsymbol p^\dagger\boldsymbol p$ and
$\chi=\boldsymbol p^T\boldsymbol p$. The oscillatory quadratic forcing
is proportional to $A^2$, rather than $|A|^2$. Its orthogonal angular
projections satisfy the exact identities
\begin{equation}
 \|P_0(A^2)\|_{L^2(S^2)}^2=\frac{|\chi|^2}{4\pi},\qquad
 \|P_2(A^2)\|_{L^2(S^2)}^2
 =\frac{3\nu^2-|\chi|^2}{10\pi}\geq\frac{\nu^2}{5\pi}.
 \label{nl:polarization-norms}
\end{equation}
These follow by decomposing $\boldsymbol p\boldsymbol p^T$ into its trace
and symmetric traceless parts and integrating the fourth spherical
moment. Circular polarization, for example
$\boldsymbol p=(1,i,0)/\sqrt2$, has $\chi=0$ and hence no oscillatory
$L=0$ source, but its $L=2$ source does not vanish.

This source identity has a conditional implication for the outgoing
response. Define $C_{\mathrm{ref}}\in\mathbb C^2$ to be the \emph{exact}
radial $L=2,m=0$ outgoing coefficient for the normalized $Y_{10}$ tangent,
with the same radial normalization. Assuming the exact regular outgoing
problem exists uniquely, spherical symmetry makes its radial operator
independent of $m$. If $b_m=\langle Y_{2m},A^2\rangle$, linearity gives
\begin{equation}
 C_{2m}=\sqrt{5\pi}\,b_m C_{\mathrm{ref}},\qquad
 \sum_m\|C_{2m}\|^2
 =\frac{3\nu^2-|\chi|^2}{2}\|C_{\mathrm{ref}}\|^2
 \geq\nu^2\|C_{\mathrm{ref}}\|^2.
 \label{nl:polarization-response}
\end{equation}
The same identity holds separately for either radial component.
Conjugating the second computational channel to the physical sideband
permutes conjugate spherical harmonics and preserves these norm sums.
Thus an exact radial $L=2$ nonzero certificate would cover every nonzero
complex polarization in this one radial dipole eigenspace. The numerical
coefficients reported below do not provide that certificate.

For clarity, the corresponding nonlinear obstruction requires a precise
continuation class. Suppose a family is twice continuously differentiable
in amplitude as a $2\pi$-periodic rotating-frame profile in
$C^2(\mathbb T;H^4_\gamma(\mathbb R^3))$, where
$H^4_\gamma=\{w:e^{\gamma\sqrt{1+|x|^2}}w\in H^4\}$.
For the exact T2 background and tangent we may fix $\gamma=1/10$:
the inherited exterior decay rates are at least $0.495$ and $0.287$,
respectively. Their radial equations propagate this decay to derivatives
through order four; the regular origin expansions are even for the
background and of the form $r^2$ times even functions for the reduced
dipole components. Consequently the unreduced Cartesian fields are
smooth at the origin and belong to this same weighted space.
This membership statement does not assume that a nonlinear continuation
exists or certify a numerical Sobolev norm. Allow smooth
carrier and modulation frequency shifts and an additional static first
tangent, but no other first-order temporal harmonics. A nonzero exact
second-harmonic outgoing amplitude would contradict localization of the
second amplitude derivative. At this order the frequency-shift terms
are static or fundamental, so they cannot remove the second-harmonic
source. This excludes only that uniformly localized smooth continuation
class; it is not a dynamical instability or lifetime theorem.

\subsection{Numerical second-harmonic response of the dipole mode}
\label{nl:dipole-response}
We evaluated the T2 forcing in both $L=0$ and $L=2$ using separately
reconstructed numerical profiles near the certified point. The reduced
real mode components obey $\int_0^\infty(u^2+v^2)\,\dd r=1$, including
an approximate tail contribution, with $\int Y_{10}^2\,\dd\Omega=1$.
The solver uses the complex conjugate of our positive-carrier convention:
both computational unknowns have positive-$q$ outgoing Hankel conditions,
and the second unknown is conjugated to reconstruct the physical negative
sideband. The power is invariant under this convention change.

Three prescribed settings use exterior radii $48,64,64$, source spacings
$0.02,0.02,0.01$, and relative ODE tolerances
$2\times10^{-10},2\times10^{-10},2\times10^{-11}$. Two further solves at
the finest settings use the coarser reconstruction level from the same
input algorithm, with different exterior radius and tolerance. This is
an input-sensitivity comparison, not an independent implementation.
All four outgoing coefficients remain nonzero under these changes;
the largest observed relative complex-amplitude change is below
$7.8\times10^{-8}$. This sensitivity is not a rigorous error bound.

At the finest setting the harmonically averaged power coefficients are
\begin{center}
\begin{tabular}{ccc}
\toprule Angular sector&Positive sideband&Negative sideband\\\midrule
$L=0$&$2.4239\times10^{-7}$&$6.0359\times10^{-5}$\\
$L=2$&$4.4956\times10^{-7}$&$9.1587\times10^{-5}$\\\bottomrule
\end{tabular}
\end{center}
The leading power is $\epsilon^4$ times the sum of these coefficients in
the stated normalization. Both physical sidebands carry positive outward
energy with the stated outgoing conditions. These values neither determine a lifetime nor measure a
nonlinear time-evolution loss. They must not be compared directly with
coefficients using the earlier radial maximum-amplitude normalization.

This calculation supports second-harmonic leakage for the approximate T2
input, beyond merely finding open channels and nonzero local sources.
It does not certify leakage at the exact enclosed BIC. Float rounding
places some input parameters outside the tiny certified box; sampled
scaled Hermite-interpolant ODE defects have a sampled maximum of about
$5.2012\times10^{-6}$ (rounded upward to $5.3\times10^{-6}$).
These are midpoint samples on $0.01$-spaced grids extending to $r=80,96$,
excluding $|r-8|\le0.011$ around the splice, not uniform residual bounds.
Differences between reconstruction levels are much smaller. Omitted exterior
source and potential contributions were diagnosed, not converted into
propagated amplitude-error bounds in that calculation. The separate
validated monopole transfer in Section~\ref{nl:conditional-monopole}
addresses these errors using a fixed comparison function and the exact
mode's Jost normalization; it does not certify the numerical power table.

A subsequent fixed-coefficient reconstruction checked the source--current
identity for each angular sector. For the real symmetric forced operator,
$J=\operatorname{Im}(U^\dagger U')$ satisfies
$J'=\operatorname{Im}(F^\dagger U)$, where $U$ denotes the two reduced
response components and $F$ their forcing. Both computational channels
use positive-wavenumber outgoing waves. Keeping the saved matching
coefficients fixed, an independent volume quadrature agrees with
$\sum_j q_j|C_j|^2$ to relative residuals $9.6\times10^{-11}$ for $L=0$
and $1.3\times10^{-11}$ for $L=2$. This is a numerical consistency check
of the finite forced problem, not an energy-power measurement or a
validated error bound. It does not independently validate a shared source
formula. The change between the two volume quadratures is
$1.9\times10^{-9}$ or smaller relative to the respective outgoing current;
this observation is not a certified quadrature-error estimate.

% Conditional enclosure incorporated in v0.8; scope remains conditional.
% Numerical error budget does not remove the inherited analytical hypotheses.
\subsection{A conditional enclosure for the monopole response}
\label{nl:conditional-monopole}
For the real dipole tangent at T2, we assembled a complete error budget
for the $L=0$ second-harmonic forced problem. Here $L$ labels the
radiated angular sector, whereas the underlying localized linear mode
has $l=1$. The calculation uses the exact-mode Jost normalization,
with the evanescent reduced component satisfying
$\lim_{r\to\infty}e^{\kappa_c r}v(r)=1$.
It therefore has a different amplitude normalization from the numerical
power table in Section~\ref{nl:dipole-response}.

The error identity includes the inner and outer source differences,
operator differences, the residual of the fixed approximate response,
its actual matching jumps and origin contribution, and the exterior
boundary and coefficient-extraction errors. The exterior potential and
source are included through infinity. No measured small linewidth or
comparison between two numerical resolutions substitutes for these terms.

The following inherited hypotheses specify the exact problem to which
the numerical error budget applies; they are not new conclusions of
the scalar summation:
\begin{enumerate}
\item The T2 existence enclosure of Section~\ref{sec:certified-existence}
supplies one common exact parameter tuple
$z_*=(a_*,c_*,\rho_*,\omega_*)$ in its certified dyadic box, a positive
regular spherical profile $f_*$, and a nonzero real regular dipole mode.
All profile, mode and parameter-error bounds refer to this same tuple.
Neither uniqueness of the tuple nor the claim that every point of the
box is an exact root is assumed.

\item At that exact tuple, the decaying fundamental Jost solution belongs
to the span of the two regular origin solutions. This is the inherited
regular/Jost identification from the T2 Wronskian conditions, not an
inference from overlap of numerical residual intervals. Its reduced
components are $O(r^2)$ at the origin and obey the same normalization
$e^{\kappa_c r}v(r)\to1$, with
$\kappa_c=\sqrt{1-(\omega_*-\rho_*)^2}$.
The background and mode satisfy the analytic infinite-tail bounds:
for $r\ge36$ one may use
$|f_*(r)|\le(97/8)e^{-0.495r}/r$ and
$|u(r)|+|v(r)|\le1.097e^{-0.287r}$.

\item The exact quadratic $L=0$ source is the complete product source
of the real standing dipole in Eqs.~\eqref{nl:quadratic}--\eqref{nl:forced}.
The fixed nominal mode, source, response and comparison coefficient
use the same amplitude convention: the nominal unit-$L^2$ mode is divided
by the exact rational
$s_0=2558188928468177/2251799813685248$, and quadratic quantities by $s_0^2$.
This does not identify $s_0$ with the norm of the exact mode.
The computational carrier conjugation is retained throughout.

\item For the second harmonic, the exact reduced operator is
$-\partial_r^2-Q_*^2+V_*(r)$ with real symmetric, locally bounded
$V_*$ and
$Q_*=\operatorname{diag}(q_+,q_-)$,
$q_\pm=\sqrt{(\omega_*\pm2\rho_*)^2-1}$.
Both channels satisfy $5/2<q_\pm<23/5$.
The exterior Volterra bounds hold at these exact frequencies,
with the analytic potential and source tails retained through infinity.
At $R=64$, the exterior potential satisfies
$\int_R^\infty\|V_*(r)\|_2\,dr/(5/2)<1/2$ and the source is integrable.
Thus the exact outgoing reduction includes both its homogeneous maps
and its inhomogeneous source corrections. Together with regular
Dirichlet data at the origin, the positive outgoing matrix current
gives the unique forced solution and adjoint extraction problem.
This argument concerns two open second-harmonic channels, not the
fundamental embedded mode.

\item The analytic and validated input bounds used in the error identity
enclose the entire core and exterior, including exact-to-nominal
frequency and map changes. In Euclidean vector and induced matrix norms,
the inherited bounds
$\int_0^\infty\|V_*(r)\|_2\,dr<25/4$ and the Volterra estimates give
the adjoint amplitude weights $\|Z\|_2<26/5$ and $\|Z'\|_2<14$.
Source and operator differences, response residuals, origin and matching
jumps, exterior source terms, and coefficient-extraction rounding are
all bounded in this same convention, without dropping a tail or counting
one twice. The directed-arithmetic implementation and its interval
libraries remain part of the computational trust assumptions.
\end{enumerate}
Under these inherited hypotheses, the computed common coefficient-error
bound is

\begin{equation}
 \|C_* - C_{\mathrm{num}}\|_2
 \leq \mathcal E < 0.001023269269845280.
\end{equation}
The frozen comparison coefficient in the second computational channel
has the lower bound
\begin{equation}
 |C_{\mathrm{num},2}| > 0.001582591541472706.
\end{equation}
The directed arithmetic establishes, before rounding for display,
\begin{equation}
 |C_{*,2}| > 0.000559322271627427.
\end{equation}
This channel has laboratory frequency $\omega_*-2\rho_*<0$.
The solver uses a negative carrier and conjugates its second computational
component to recover that physical sideband. Conjugating the entire field
to the positive-carrier convention used here gives
$C_{-,\mathrm{paper}}=C_{*,2}$; all these operations preserve the modulus.
The first channel remains undecided by this common error bound; that
failure does not imply its coefficient vanishes.

This result supplies the nonzero-coefficient premise of the conditional
smooth-continuation argument for a real standing dipole. Its extension
to a complex dipole polarization is proportional to
$\boldsymbol p^T\boldsymbol p$ and therefore does not cover circular
polarization. It gives neither a nonlinear lifetime nor a dynamical
instability theorem. In particular, the bound is not a certified value
for any power coefficient in the preceding approximate-input table.


\subsection{Existing evidence for radial location C0}
\label{nl:evidence}
Separate computations already address the earlier $n=1$, $l=0$ mode
near $\omega^2=0.7976767871$ and $\rho=1.7446175448$. They concern the
same sextic field equation, not T1 or T2. A finite-amplitude radial
evolution compared the nonlinear field with its background and linear
response on two spatial grids. For amplitudes $0.01,0.02,0.04$, the
integrated quadratic flux of the nonlinear remainder increased by
factors approximately $15.97$ and $15.89$ upon doubling amplitude.
This supports fourth-order scaling over the tested range; that remainder
flux is not an exact measurement of the total lost core energy.

A subsequent harmonic diagnostic identified the dominant laboratory
sideband $E=\omega-2\rho\simeq-2.59611$ as outgoing at two detector
radii, on both grids and in both prescribed time windows. The weaker
$E=\omega+2\rho$ component was resolved but did not satisfy every
outgoing-wave criterion. Large residuals in the full field's harmonic
fit remained, so these observations do not establish a periodic
multifrequency solution.

A separately solved frequency-domain problem, using the same background
and mode data but no fitted amplitude from the time evolution, predicted
the dominant outgoing coefficient. At $\epsilon=0.02$, its predicted
amplitude was $5.1563\times10^{-5}$, compared with measured values
approximately $(5.03$--$5.10)\times10^{-5}$. The prescribed comparison
criteria passed. Shared input profiles make this a separate mechanism
test, not a wholly independent reconstruction of the background.

An interval calculation further separated one necessary overlap from
zero for an explicitly frozen finite operator on $[0,44]$:
\begin{equation}
 I_2\in[-0.16383470526432667,-0.10547329373338812].
 \label{nl:frozen}
\end{equation}
The corresponding $I_1$ interval contained zero, and the original test
requiring both exclusions remained failed. Nevertheless, one certified
nonzero overlap suffices to violate the necessary simultaneous condition
$I_1=I_2=0$ for that frozen problem. Its binary64 spline inputs are treated
as exact data, with a free exterior and zero source beyond radius $44$.
This is not a certificate that the exact continuum Q-ball and its BIC
have that overlap: a validated link to the continuum inputs and exterior
tail remains required. The linear existence certificates do not, by
themselves, supply that quantitative link.

If a leading second-order outgoing amplitude is nonzero, its power is
of order $\epsilon^4$. A lifetime law additionally requires a controlled
relation to mode energy and slow amplitude evolution, and separation of
initial adjustment from persistent radiation. Neither the finite-time
observations nor the frozen-operator certificate provides a general
nonlinear stability or lifetime theorem. Earlier real-field shell-model
normal forms describe a different model and supply no substitute for
these Q-ball source and overlap conditions.
